% --- Mode d'emploi ----------------------------------------------------------
\chapter*{Comment ce livre est fait}
\markboth{Comment ce livre est fait}{Comment ce livre est fait}
\addcontentsline{toc}{chapter}{Comment ce livre est fait}

Tous les chapitres suivent la même route, dans le même ordre. Au bout de deux
ou trois chapitres, vous saurez où regarder sans chercher.

\vspace{0.8em}

\begin{revision}
On rouvre ce que vous savez déjà et dont le chapitre va se servir : une
formule de première, une propriété de seconde. Jamais un mot de la notion
nouvelle.
\end{revision}

\begin{automatismes}
Une dizaine de calculs qui doivent sortir sans papier. C'est le chronomètre,
pas la réflexion, qui compte ici.
\end{automatismes}

\noindent Vient ensuite \textbf{le cours}, où alternent plusieurs formes :

\begin{definition}[Ce qu'on pose]
Le mot nouveau, écrit une fois pour toutes. On y revient chaque fois qu'un
doute apparaît.
\end{definition}

\begin{theoreme}[Ce qu'on démontre]
Le résultat qui fait travailler la définition, suivi de sa preuve.
\end{theoreme}

\begin{demonstration}
La preuve est écrite en petit, en retrait. Elle se lit après le théorème, ou
plus tard, ou jamais si le temps manque — mais elle est là.
\end{demonstration}

\begin{visualisation}[la figure qui montre l'idée]
Une figure, et à côté la phrase qu'un professeur dirait en la traçant au
tableau. Regardez d'abord le dessin, lisez la légende ensuite.
\end{visualisation}

\begin{exemple}[le calcul déroulé en entier]
Un cas traité ligne à ligne, sans étape sautée.
\end{exemple}

\begin{methode}{reconnaître le bon outil}
Quand plusieurs résultats s'appliquent, ce bloc dit lequel choisir et à quoi
on le reconnaît. Ce sont ces pages-là que l'on relit la veille de l'épreuve.
\end{methode}

\begin{erreurclassique}
La faute que l'on corrige chaque année sur les copies, écrite noir sur blanc
pour que vous la voyiez venir.
\end{erreurclassique}

\noindent Puis on passe la main. Six formats d'exercices se succèdent, et ce
n'est pas pour faire joli : chacun teste une chose que les autres ne testent
pas.

\begin{entrainement}[juste après la notion]
Des exercices courts, placés immédiatement après la notion qu'ils font
travailler, classés du plus simple au plus exigeant. Les pastilles
\niveau{1}, \niveau{2}, \niveau{3} donnent le niveau.
\end{entrainement}

\begin{qcm}
Quatre propositions, une seule juste. On y perd du temps si l'on calcule
tout ; on en gagne si l'on sait éliminer. C'est un entraînement au tri, pas
au calcul.
\end{qcm}

\begin{vraifaux}
Un énoncé, et il faut trancher — puis dire pourquoi en une ligne. Un contre-%
exemple suffit à réfuter ; une impression ne suffit jamais à confirmer.
\end{vraifaux}

\begin{lecturegraphique}
Une courbe est donnée, et tout se lit dessus : un ensemble de définition, un
signe, une limite, une tangente. Aucun calcul n'est demandé. L'épreuve
commence souvent par là.
\end{lecturegraphique}

\begin{situation}
Un énoncé venu d'ailleurs que du cours : une récolte, un effectif scolaire,
un prix qui monte. Les chiffres sont plausibles, la question est utile, et la
mise en équation vous appartient.
\end{situation}

\begin{annales}
Des exercices tirés de sujets réellement posés. Rien n'y est inventé : chaque
énoncé porte sa session et son numéro. C'est le meilleur thermomètre dont
vous disposiez.
\end{annales}

\begin{jecherche}[une situation à démêler]
Un problème plus ouvert, où plusieurs chemins mènent au résultat, et où l'on
ne vous dit pas lequel prendre.
\end{jecherche}

\begin{jemeteste}
Un questionnaire de fin de chapitre, à faire seul, pour savoir si le chapitre
est acquis avant de passer au suivant.
\end{jemeteste}

\begin{commeaubac}
Un exercice au format et au barème de l'épreuve.
\end{commeaubac}

\begin{devoirsurveille}
À la fin de chacune des quatre parties, un devoir complet : deux heures
quinze, deux exercices de cinq points, un problème de dix, exactement comme
le jour de l'examen. Sa portée est cumulative — il peut mobiliser tout ce qui
précède dans le livre, jamais ce qui suit.
\end{devoirsurveille}

\begin{notepourlaclasse}
Ce bloc n'apparaît que dans l'édition du professeur : durée conseillée,
difficulté prévisible, variante pour classe à grand effectif.
\end{notepourlaclasse}

\vspace{0.6em}

\begin{remarque}[ce que ce livre ne contient pas]
Pas de bloc « je vise le concours » : aucun sujet d'entrée à l'ENI ou à
l'ESPA ne correspond au profil de la série littéraire. Pas de nombres
complexes, pas de géométrie dans l'espace, pas de calcul matriciel, pas de
probabilité conditionnelle ni de variable aléatoire : le programme officiel
les déclare hors programme pour cette série, et aucun sujet ne les a jamais
demandés.
\end{remarque}
